\SetPractica{Grupos funcionales orgánicos}{Identificación de grupos funcionales orgánicos}

\section{Introducción teórica}



\section{Objetivos}

\begin{itemize}
    \item Familiarizar al alumno con algunas de las reacciones características que permiten la identificación de ciertos grupos funcionales orgánicos 
    \item Aplicar este conocimiento para aprender a identificar el tipo de compuesto presente en una serie de muestras
\end{itemize}

\section{Materiales, equipos y reactivos}

\begin{table}[H]
\centering{\small\setstretch{1.25}
\begin{tabular}{|m{0.2\linewidth}|m{0.2\linewidth}|m{0.45\linewidth}|} \hline
    
    \thead{Equipos} & \thead{Materiales} & \thead{Reactivos} \\ \hline
    
    {Balanza \par}
    {Campana de extracción de gases \par}
    {Termobaño}
    
    &
    
    {Probetas \par}
    {Vasos de precipitado \par}
    {Tubos de ensayo \par}
    {Pipetas Pasteur \par}
    {Pipetas graduadas \par}
    {Micropipetas 1mL \par}
    {Pinzas de moss o para tubos de ensayo \par}
    {Espátulas}
    
    & 
    
    {Cloruro de acetilo (\ce{H3CCOCl}) \par}
    {Disolución de almidón 0.1\% \par}
    {Disolución de cloruro de hierro III (\ce{FeCl3}), al 10\% \par}
    {Disolución de 2,4-dinitrofenilhidracina (DNPH) \par}
    {Disolución de hidróxido amónico (\ce{NH4OH}) 2M \par}
    {Disolución de hidrocloruro de hidroxilamina (\ce{NH3OHCl}) 1N \par}
    {Disolución de ácido clorhídrico (\ce{HCl}) 2N \par}
    {Disolución de ioduro de potasio (\ce{KI}) al 2\% \par}
    {Disolución de iodato de potasio (\ce{KIO3}) al 4\% \par}
    {Disolución de hidróxido de potasio (\ce{KOH}) en metanol 2N \par}
    {Disolución de hidróxido de sodio (\ce{NaOH})  al 5\%  \par}
    {Disolución de hidróxido de sodio (\ce{NaOH})  al 10\% \par}
    {Disolución de nitrato de plata   (\ce{AgNO3}) al 5\%  \par}
    {Disolución de nitrato de plata   (\ce{AgNO3}) al 10\% \par}
    {Metanol \par}
    {Sulfato amónico de hierro II (\ce{Fe(SO4)2(NH4)2}) \par}
    {Reactivo de Lucas \par} 
    {Reactivo de Brady} 
    
    \\ \hline
\end{tabular}}\end{table}


\section{Procedimientos}

\subsection{Preparación previa de reactivos importantes}

\subsubsection{Preparación del reactivo de Lucas}
\begin{enumerate}
    \item Se añaden 68g de \ce{ZnCl2}, en porciones a 45ml de HCl concentrado, procurando enfriar la disolución a lo largo del proceso.
\end{enumerate}

\subsubsection{Preparación del reactivo de Brady}
\begin{enumerate}
    \item Disolver 1g de DNPH (2,4- dinitrofenilhidrazina) en 5ml de \ce{H2SO4} concentrado
    \item Se añaden lentamente y con agitación esta mezcla a una disolución de agua (7ml) y etanol (25ml)
    \item Finalmente se filtra para eliminar sólidos en suspensión.
\end{enumerate}

\subsubsection{Preparación del reactivo de Tollens}
\begin{enumerate}
    \item En un tubo pequeño, añadir 2 gotas de una disolución de \ce{NaOH} al 5\% a 1ml de disolución acuosa de \ce{AgNO3} al 5\%.
    \item Agitar el tubo y añadir gota a gota y con agitación \ce{NH4OH} 2N, hasta que se consiga disolver el precipitado de AgOH, que previamente se había formado.
\end{enumerate}



\subsection{Detección de alcoholes: Ensayo con cloruro de acetilo}
\begin{enumerate}
    \item En un tubo de ensayo pequeño seco, colocar la muestra problema (aprox. 0,5mL de líquido)
    \item Añadir 0,3mL cloruro de acetilo, gota a gota.
    \begin{enumerate}
        \item[!] \textcolor{red}{\textbf{ADVERTENCIA:}} La adición de cloruro de acetilo debe hacerse dentro de la campana de extracción de gases y empleando guantes.
    \end{enumerate}
    \item La aparición de reacción, indicada por un ligero calentamiento de la mezcla y por la acidificación de la misma, es indicación positiva para la presencia de un alcohol.
\end{enumerate}

\subsubsection{Diferenciación entre alcoholes primarios, secundarios y terciarios}
\begin{enumerate}
    \item Tomar 0,5mL (o aprox. 500mg) de la muestra problema en un tubo de ensayo pequeño
    \item Añadir 3mL del reactivo de Lucas
    \begin{enumerate}
        \item[!] \textcolor{red}{\textbf{PRECAUCIÓN:}} Contiene \ce{HCl} concentrado.
    \end{enumerate}
    \item Cerrar el tubo con un tapón y agitar durante unos 15 segundos.
\end{enumerate}

\begin{table}[H]
\centering{\small\setstretch{1.5}
\begin{tabular}{|m{0.4\linewidth}|m{0.4\linewidth}|}
    \thead{Si la solución se enturbia rápidamente} & \thead{Si la solución permanece clara} \\ \hline
    Alcohol terciario & Alcohol primario o secundario \\ \hline
\end{tabular}}
\end{table}

\textbf{Nota:} Debe indicarse que este ensayo, aunque puede utilizarse para detectar alcoholes secundarios, es poco fiable. Además, no es válido para alcoholes acíclicos ni para los insolubles en agua.




\subsection{Detección de aldehídos y cetonas: Formación de 2,4-dinitrofenilhidrazonas (test de Brady)}

\subsubsection{Test de Brady}
\begin{enumerate}
    \item En un tubo de ensayo pequeño, disolver 2 ó 3 gotas de la muestra problema (o aprox. 50mg de sólido) en 0,5mL de metanol.
\item Añadir 1mL del reactivo de Brady y agitar vigorosamente.
\item La aparición de un precipitado amarillo, naranja o rojo indica la presencia de un compuesto carbonílico.
\end{enumerate}

\subsubsection{Diferenciación entre aldehídos y cetonas (test de Tollens)}
\begin{enumerate}
    \item Limpiar el tubo de ensayo pequeño donde se realizará el experimento, hirviendo en éste \ce{NaOH} al 10\%.
    \item Añadir 3 gotas de la muestra problema (o bien 50mg de sólido), y 2.5mL de reactivo de Tollens recién preparado.
    \item Agitar el tubo y dejarlo reposar durante 10 minutos.
    \item Si al cabo de este tiempo no se observa reacción, se calienta el tubo en un baño de agua caliente.
    \item La formación de un precipitado negro, o de un espejo de plata en las paredes del tubo de ensayo, constituye una prueba positiva de la presencia de un aldehído.
    \item Es importante lavar los tubos utilizados con \ce{HNO3} diluido y desechar los restos sobrantes de la disolución.
\end{enumerate}

\bigskip

\textcolor{red}{\textbf{PRECAUCIÓN:}} El isocianato de plata, compuesto muy explosivo en seco, puede estar presente en los residuos de las soluciones del reactivo de Tollens; por esta razón, es conveniente tirar el resto de la disolución de Tollens no utilizada y lavar los tubos con \ce{HNO3} diluido.




\subsection{Detección de ácidos: Ensayo del yodato-yoduro para ácidos carboxílicos}
\begin{enumerate}
    \item Tomar 2 gotas ó 5mg de la sustancia problema y colocarla en un tubo de ensayo pequeño.
    \item Añadir 2 gotas de una solución de KI al 2\% y 2 gotas de una solución de KIO3 al 4\% (ambas en \ce{H2O}).
    \item Se tapa el tubo y se calienta en un baño de agua hirviendo durante 1 minuto.
    \item Enfriar el tubo y añadir 2 gotas de solución de almidón, recientemente preparada, al 0,1\%.
    \begin{enumerate}
        \item Si la sustancia es un ácido carboxílico, aparecerá un color azul.
    \end{enumerate}
\end{enumerate}




\subsection{Detección de ésteres: Ensayo para la detección de ésteres}
\begin{enumerate}
    \item Colocar 0,5mL de solución de clorhidrato de hidroxilamina 1N en etanol en un tubo de ensayo.
    \item Adicionar 30mg (2-3 gotas) del compuesto que va a ser ensayado.
    \item Adicionar, gota a gota, una solución de KOH 2N en metanol, hasta que la mezcla tenga pH 9-10
    \item Una vez alcanzado el pH anteriormente mencionado, adicionar 4 gotas más de KOH 2N.
    \item Calentar la mezcla hasta la ebullición (en un termobaño).
    \item Enfriar y añadir gota a gota y con agitación HCl 2N hasta llegar a un pH aproximado de 3.
    \item Añadir 1 gota de solución de FeCl3 al 10\% y observar la coloración.
    \begin{enumerate}
        \item Una coloración rojo sangre es indicativo de la presencia de ésteres.
    \end{enumerate}
\end{enumerate}